% 【本文件写什么】api-v0 契约层，叶子，只写语义不写实现
% - crate 路径：os/components/wateros-cred/cred-api/api-v0/src/
% - 列出并说明：pub trait、核心类型、错误枚举、常量
% - 每个公开 API 的前置条件、后置条件、线程/中断上下文限制
% - 与 Linux/用户态约定的对齐点与 deliberate 差异
% - v0 版本当前行为与后续可替换 extension 点
% 【事实来源】
% - 上述 crate 下全部 .rs 源文件
% - docs/exports/public-api/ 中对应符号表，以源码为准
% 【不写什么】
% - impl 内部数据结构、汇编、具体算法步骤

\paragraph{基础类型}

\code{Uid(u32)}、\code{Gid(u32)}对齐Linux 32位ID语义，与\code{stat}/syscall返回一致。\code{TaskId = usize}为cred侧表索引。\code{SUPPLEMENTARY\_GROUP\_COUNT = 32}限制supplementary组数量。

\paragraph{ProcessCredentials}

快照包含八ID，real/effective/saved/fs的uid与gid及最多32个supplementary组。\code{ROOT}常量：全部ID为0，\code{supplementary\_group\_len = 1}且组为\code{[0]}。

实例方法实现privileged语义的ID更新，与Linux \code{set*id} 对齐：
\begin{itemize}
 \item \code{set\_uid}/\code{set\_gid}：四处ID一并更新；
 \item \code{set\_reuid}/\code{set\_regid}：\code{None}表示Linux \code{-1}保持不变，并同步saved/fs；
 \item \code{set\_resuid}/\code{set\_resgid}：含saved三元组；
 \item \code{set\_supplementary\_groups}：替换组列表；
 \item \code{supplementary\_group\_count()}：供\syscall{getgroups(0)}返回计数。
\end{itemize}

\paragraph{Capability}

枚举\code{Capability::Chown}、\code{Capability::SysAdmin}，与\code{capget}/\code{capset}及xattr管理操作预留对齐。\code{prctl(PR\_CAPBSET\_*)}权威语义最终将归口\code{AccessCheck::has\_cap}。

\paragraph{CredentialBackend}

\code{current(tid)}读取指定任务快照；无条目时由impl决定，\code{impl-root} 为 \code{panic}。生命周期：\code{on\_user\_task\_spawned}、\code{fork\_cred}、\code{on\_exec}、\code{drop\_task\_cred}。须在任务上下文或明确的spawn/reap路径调用；不可从中断上下文无保护地变更侧表。

\paragraph{CredentialMutation}

per-task的\code{set\_resuid}、\code{set\_resgid}、\code{set\_supplementary\_groups}；\code{None}参数语义同Linux \code{-1}。聚合层\code{set\_uid}等门面在调用前解析当前\code{TaskId}。

\paragraph{AccessCheck}

\code{has\_cap}、\code{may\_access\_inode}、\code{may\_chown}为权限查询扩展点。\code{impl-root}阶段：\code{has\_cap}对effective uid为0返回true，否则false；\code{may\_access\_inode}恒true；\code{may\_chown}对root/\code{CAP\_CHOWN}放行，非root仅允许有限属主变更。真实capability位图与VFS open权限校验为后续工作。

\paragraph{版本与差异}

v0对齐Linux cred子集，不含namespace、securebits、完整cap位图。\code{on\_exec}契约保留S\_ISUID/S\_ISGID解析扩展点，即 \code{TODO(cred-exec-setuid)}，当前impl为no-op。
